% ==============================================================================
% فصل اول: مقدمه و بیان مسئله
% ==============================================================================

\chapter{مقدمه و بیان مسئله}
\label{ch:intro}

\section{مقدمه و ضرورت پژوهش}
در دهه‌های اخیر، سامانه‌های تعبیه‌شده (\lr{Embedded Systems}) به ستون فقرات زیرساخت‌های حیاتی دنیای مدرن بدل شده‌اند؛ از سیستم‌های اویونیک و هدایت ناوبری در صنایع هوافضا و خودروهای هوشمند خودران گرفته تا تجهیزات کنترل علائم حیاتی در پزشکی و بازوهای رباتیک در خطوط اتوماسیون صنعتی. ویژگی مشترک و وجه تمایز بنیادین این رده از کاربردها، ماهیت «بی‌درنگ سخت» (\lr{Hard Real-Time}) آن‌ها است. در سامانه‌های بی‌درنگ سخت، درستی و صحت کل سامانه صرفاً به درستی محاسبات منطقی و ریاضی محدود نمی‌شود، بلکه عامل «زمان» به عنوان یک بُعد تعیین‌کننده و غیرقابل‌مذاکره وارد عمل می‌گردد. به بیان دقیق‌تر، تولید یک پاسخ کاملاً صحیح از نظر منطقی، چنانچه با حتی کسری از میلی‌ثانیه تأخیر پس از موعد مقرر (\lr{Deadline}) تحویل داده شود، در حکم خطای بحرانی سامانه بوده و می‌تواند فجایع جبران‌ناپذیر انسانی یا خسارات سنگین مالی به همراه داشته باشد.

برای تحقق ضمانت‌های قطعی در زمان‌بندی‌پذیری (\lr{Schedulability})، مهندسان سیستم‌های بی‌درنگ نیازمند دانستن کران بالای زمان اجرای تک‌تک وظایف محاسباتی هستند. این کران بیشینه، در متون علمی تحت عنوان «بیشینه زمان اجرا» (\lr{Worst-Case Execution Time - WCET}) شناخته می‌شود \cite{wilhelm2008worst}. استخراج دقیق و قابل‌اطمینان \lr{WCET} سنگ‌بنای طراحی سامانه‌های با قابلیت اطمینان بالا (\lr{High-Reliability}) و تخصیص منابع پردازشی است \cite{kargahi2008analytical}. با این وجود، به دلیل جهش‌های فنی گسترده در معماری پردازنده‌های تعبیه‌شده در سال‌های اخیر، محاسبه و استنتاج این کمیت به یکی از پیچیده‌ترین چالش‌های مهندسی الکترونیک و کامپیوتر مبدل گشته است.

در سطح صنایع پیشرفته جهانی، خودکارسازی فرآیند آزمون محصول و استنتاج غیرتهاجمی زمان‌بندی توسط پلتفرم‌های نرم‌افزاری انحصاری و بسیار پرهزینه نظیر سامانه‌های شرکت \lr{Rapita Systems} (مجموعه \lr{RVS / RapiTime})، بسترهای آزمون در حلقه سخت‌افزار (\lr{HIL}) شرکت \lr{Speedgoat} و ابزارهای مشابه انجام می‌پذیرد. هزینه لایسنس و به‌کارگیری این راهکارهای تجاری بالغ بر ده‌ها تا صدها هزار دلار (بیش از ۱۰ تا ۱۰۰ هزار دلار) بوده و دستیابی به آن‌ها نیازمند بودجه‌های سنگین است. از این رو، طراحی و توسعه یک پلتفرم جامع آزمون خودکار محصول و پایش غیرتهاجمی با حداقل هزینه ممکن، اهمیتی دوچندان و توجیه اقتصادی و کاربردی بسیار بالایی پیدا می‌کند.

\section{مسئله استنتاج \lr{WCET} و چالش‌های معماری پردازنده‌های مدرن}
در نسل‌های نخستین میکروکنترلرها (نظیر پردازنده‌های ۸ بیتی و ریزکنترلرهای اولیه)، خط سیر اجرای دستورات به شدت تعیینی (\lr{Deterministic}) بود؛ به طوری که هر دستورالعمل تعداد چرخه‌های کلاک مشخص و ثابتی را مصرف می‌کرد و طراح سیستم می‌توانست با جمع جبری زمان دستورات و بازرسی چشمی کد اسمبلی، بدترین زمان اجرای یک تابع را با دقت بالا تخمین بزند. اما در پردازنده‌های امروزی تعبیه‌شده با کارایی بالا (نظیر هسته‌های \lr{ARM Cortex-M7})، مجموعه‌ای از نوآوری‌های سخت‌افزاری جهت افزایش توان محاسباتی به کار گرفته شده است که رفتار زمانی پردازنده را به شدت پویا، غیرخطی و وابسته به کانتکست تاریخی اجرا می‌نماید:
\begin{enumerate}
    \item \textbf{خط‌لوله‌های چندمرحله‌ای و سوپراسکالر (\lr{Superscalar Pipelines}):} اجرای هم‌زمان بیش از یک دستور در چرخه کلاک و وابستگی‌های ساختاری داده‌ها (\lr{Data Hazards}).
    \item \textbf{حافظه‌های نهان چندسطحی (\lr{Instruction \& Data Caches}):} وابستگی زمان واکشی دستور به رخداد اصابت (\lr{Cache Hit}) یا عدم اصابت (\lr{Cache Miss})، که موجب نوسان زمان دسترسی از چند نانوثانیه تا ده‌ها چرخه کلاک می‌گردد.
    \item \textbf{واحدهای پیش‌بینی انشعاب (\lr{Branch Predictors}):} حدس زدن نتیجه شرط‌ها پیش از حل شدن آن‌ها و نیاز به پاک‌سازی خط‌لوله (\lr{Pipeline Flush}) در صورت پیش‌بینی نادرست.
    \item \textbf{تداخل گذرگاه‌ها و ساختارهای دسترسی مستقیم به حافظه (\lr{DMA}):} رقابت بر سر پهنای باند حافظه مشترک میان هسته و تجهیزات جانبی.
\end{enumerate}

این پدیده‌ها مجموعاً موجب بروز اثراتی موسوم به «ناهنجاری‌های زمانی» (\lr{Timing Anomalies}) می‌گردند؛ وضعیتی که در آن یک رخداد با فرض بهینه‌سازی محلی (مانند اصابت به حافظه نهان)، در مقیاس کلی به دلیل تغییر ترتیب تخصیص منابع منجر به تأخیر در اجرای کلی برنامه می‌شود \cite{cullmann2010predictability}. در چنین محیط‌های پویایی، استخراج کران بالای زمان اجرا صرفاً با استفاده از مدل‌های نظری ریاضی، یا به ارقام غیرواقعی و بیش‌ازحد محافظه‌کارانه (\lr{Over-conservative}) می‌انجامد که موجب اتلاف عظیم سخت‌افزار می‌شود، و یا به دلیل خطاهای تقریب مدل سخت‌افزاری، کران نامعتبری تولید می‌کند که نقض‌کننده ایمنی سامانه است.

\section{معضل ابزاردقیق نرم‌افزاری و پدیده «اثر پروب»}
برای رفع محدودیت‌های روش‌های صرفاً تحلیلی، مهندسان به روش‌های مبتنی بر اندازه‌گیری تجربی (\lr{Measurement-Based Timing Analysis - MBTA}) روی آورده‌اند \cite{bernat2002wcet}. در این رویکرد، کد برنامه بر روی بستر سخت‌افزاری واقعی بارگذاری شده و با ورودی‌های گوناگون تحریک می‌گردد تا زمان‌های اجرای آن مستقیماً ثبت شود.

با این حال، بزرگ‌ترین مانع در پیاده‌سازی این روش، نحوه پایش و ثبت زمان‌ها و شناسایی مسیر پیموده‌شده است. روش سنتی مهندسی نرم‌افزار، استفاده از «ابزاردقیق نرم‌افزاری» (\lr{Software Instrumentation}) است که در آن دستورات اضافه‌ای به سورس‌کد تزریق می‌شود؛ نظیر تغییر وضعیت پین‌های ورودی/خروجی عمومی (\lr{GPIO Toggle}) در شروع و پایان هر تابع یا ذخیره مقادیر تایمرها در آرایه‌های حافظه. این دستکاری به شدت رفتار زمانی را آلوده می‌کند. این معضل مهندسی در متون علمی با عنوان \textbf{اثر پروب (\lr{Probe Effect})} شناخته می‌شود:
\begin{itemize}
    \item افزودن دستورات پروب، خط‌لوله دستورات را مخدوش می‌کند.
    \item کدهای پروب ساختار محلی حافظه نهان (\lr{Cache Locality}) را به هم ریخته و باعث بیرون رانده شدن ناخواسته خطوط کش برنامه اصلی می‌شوند.
    \item تأخیر زمانی تحمیل‌شده توسط این ابزارها ممکن است توالی دقیق وقفه‌ها و تداخل‌های بی‌درنگ را تغییر داده و رفتاری متناقض با عملکرد واقعی برنامه در محیط عملیاتی به نمایش بگذارد.
\end{itemize}

بنابراین، برای استنتاج علمی، دقیق و قابل‌اتکای \lr{WCET}، وجود یک سازوکار پایش \textbf{کاملاً غیرتهاجمی (\lr{Non-intrusive})} که بدون تغییر حتی یک بایت از کدهای اجرایی و بدون ایجاد کوچک‌ترین تغییر در رفتار چرخه‌ای پردازنده، جریان دقیق اجرای دستورات را استخراج نماید، ضرورتی حیاتی به شمار می‌رود.

\section{معماری کلان پژوهش و مدل همکاری دوگانه}
پژوهش حاضر در قالب یک معماری یکپارچه و در تعامل با یک سامانه آزمون خودکار مکمل طراحی شده است تا چرخه کامل استنتاج \lr{WCET} را بر روی سخت‌افزار واقعی شکل دهد. این بستر جامع از دو پروژه مستقل و هماهنگ تشکیل شده است:

\begin{figure}[ht]
\centering
\begin{latin}
\begin{tikzpicture}[
    node distance=2.3cm and 1.8cm,
    block/.style={rectangle, draw=blue!80!black, fill=blue!5, thick, rounded corners=6pt, text width=6.6cm, align=center, minimum height=2.4cm, inner sep=6pt},
    testblock/.style={rectangle, draw=orange!80!black, fill=orange!5, thick, rounded corners=6pt, text width=6.6cm, align=center, minimum height=2.4cm, inner sep=6pt},
    hwblock/.style={rectangle, draw=green!60!black, fill=green!5, thick, rounded corners=6pt, text width=15.0cm, align=center, minimum height=2.2cm, inner sep=6pt},
    line/.style={draw=black!75, -latex, very thick}
]
    % Blocks
    \node [block] (provider) {
        {\large\bfseries\rl{پروژه سیستم تأمین داده}}\\[0.15cm]
        \small\rl{(موضوع این پایان‌نامه)}\\[0.2cm]
        \small\rl{استخراج غیرتهاجمی ردگیری با سخت‌افزار} \lr{ETMv4}\\[0.08cm]
        \small\rl{پردازش سیگنال، فریم‌بندی و بازسازی با} \lr{OpenCSD}
    };
    
    \node [testblock, right=1.8cm of provider] (tester) {
        {\large\bfseries\rl{پروژه سامانه آزمون خودکار}}\\[0.15cm]
        \small\rl{(پروژه مکمل)}\\[0.2cm]
        \small\rl{تولید هدایت‌شده ورودی‌ها بر پایه روش‌های ابتکاری}\\[0.08cm]
        \small\rl{ماک‌سازی نرم‌افزاری پریفرال‌ها و تزریق داده}
    };
    
    \node [hwblock, below=2.2cm of $(provider.south)!0.5!(tester.south)$] (target) {
        {\bfseries\rl{سخت‌افزار هدف: میکروکنترلر صنعتی} \lr{ARM Cortex-M7 (STM32H750)}}\\[0.2cm]
        \small\rl{اجرای کد اصلی برنامه و انتشار سیگنال‌های فیزیکی ردیابی بدون سربار و بدون وقفه}\\[0.08cm]
        \small\rl{خطوط درگاه موازي:} \lr{TRACE\_CLK} \rl{و} \lr{TRACE\_DATA[0:3]}
    };

    % Paths
    \path [line] (tester.south) -- node [right=0.15cm, align=left, font=\small] {\rl{تزریق ورودی از طریق}\\[-0.05cm]\lr{CAN / USB Mocking}} (tester.south |- target.north);
    \path [line] (provider.south |- target.north) -- node [left=0.15cm, align=right, font=\small] {\rl{ردگیری سخت‌افزاری ۵ سیمه}\\[-0.05cm](\lr{ETM Interface})} (provider.south);
    \path [line, dashed] (provider.east) -- node [above=0.12cm, align=center, font=\scriptsize] {\rl{فیدبک ردگیری}\\\rl{(پوشش و زمان)}} (tester.west);
\end{tikzpicture}
\end{latin}
\caption{مدل مفهومی معماری یکپارچه دوگانه سامانه استنتاج \lr{WCET} بر بستر سخت‌افزار واقعی}
\label{fig:dual_system_architecture}
\end{figure}

این دو پروژه عبارتند از:
\begin{enumerate}
    \item \textbf{پروژه سیستم تأمین داده (موضوع پایان‌نامه حاضر):} طراحی و استقرار \textbf{«سامانه تأمین داده‌های ردیابی دقیق و غیرتهاجمی (\lr{Data Provider System})»}. رسالت این بخش، تسخیر امواج فیزیکی ردگیری از پایه‌های میکروکنترلر، نمونه‌برداری با فرکانس بالا، بازسازی جفت‌نیبل‌ها، فیلتر بسته‌ها، حل چالش‌های همگام‌سازی و تغذیه خودکار موتور تجزیه و تحلیل \lr{OpenCSD} جهت استخراج دنباله دقیق دستورات اجرا شده و برچسب‌های زمانی با خطای رمزگشایی صفر و جیتر صفر است.
    \item \textbf{پروژه سامانه آزمون خودکار (پروژه مکمل):} طراحی و پیاده‌سازی \textbf{«سامانه آزمون خودکار جهت استنتاج \lr{WCET} (\lr{Automated Testing System})»}. این بخش به دنبال هدایت اجرای برنامه به سوی بدترین و طولانی‌ترین مسیرهای منطقی بر اساس الگوریتم‌های هوشمند جستجو و تحلیل خروجی‌های زمان‌بندی تأمین‌شده توسط پروژه اول است.
\end{enumerate}

به منظور شفاف‌سازی کامل سهم نوآوری‌های نگارنده در این کار مشترک و تبیین مرزهای ارزیابی در جلسه دفاع، جدول \ref{tab:task_division} تفکیک دقیق وظایف و خروجی‌های هر یک از این دو پروژه را نشان می‌دهد.

\begin{table}[ht]
\centering
\caption{تفکیک وظایف، مرزبندی معماری و سهم نوآوری‌های نگارنده در مقایسه با پروژه مکمل}
\label{tab:task_division}
\small
\begin{tabular}{|p{3.2cm}|p{6.8cm}|p{4.5cm}|}
\hline
\textbf{مؤلفه و حوزه فنی} & \textbf{سهم و نوآوری این رساله (پروژه تأمین داده)} & \textbf{سهم پروژه مکمل (سامانه آزمون)} \\
\hline
\textbf{سخت‌افزار و ردیابی بی‌درنگ} & 
مهندسی معکوس و راه‌اندازی ثبّات‌های \lr{ETMv4} و \lr{TPIU} در \lr{STM32H7}، حل مغایرت‌های راهنمای مرجع، تثبیت مسیرکشی پایه‌ها و نمونه‌برداری $10\times$ با لاجیک‌آنالایزر & 
تعامل با درگاه‌های پروگرامینگ جهت تزریق مقادیر اولیه و پیکربندی سناریوهای آزمون \\
\hline
\textbf{پردازش سیگنال و داده‌های خام} & 
توسعه اسکریپت تخصصی \lr{\texttt{TraceStreamProcessor.py}}، آشکارسازی لبه‌ها، فیلترینگ گلیچ، تراز فریم‌های \lr{TPIU} با الگوهای \lr{FSYNC} و تولید اسنپ‌شات استاندارد & 
الگوریتم‌های فازینگ و تولید هوشمند داده‌های ورودی برای پیمایش مسیرهای طولانی \\
\hline
\textbf{بازسازی رفتار برنامه و استخراج زمان} & 
یکپارچه‌سازی با کتابخانه مرجع \lr{OpenCSD}، انطباق با باینری \lr{ELF}، بازسازی جریان دستورات و استخراج برچسب‌های زمانی نانوثانیه‌ای و شمارش چرخه‌ها (\lr{CCI}) & 
دریافت ماتریس داده‌های ردیابی و به‌کارگیری مدل‌های تحلیلی جهت استنتاج ریاضی \lr{WCET} \\
\hline
\textbf{ابزار توسعه و اتوماسیون (\lr{IDE})} & 
طراحی و پیاده‌سازی افزونه اختصاصی \lr{CoreSight Trace Studio} در محیط \lr{VS Code} جهت تزریق خودکار درایور و پنجره‌ها و مصورسازی نتایج & 
توسعه فریم‌ورک ماک‌سازی توابع در لایه میانی \lr{HAL} جهت ایزوله‌سازی آزمون \\
\hline
\end{tabular}
\end{table}

\section{اهداف تفصیلی پژوهش}
با تکیه بر چارچوب کلان مسئله، اهداف مشخص این پایان‌نامه عبارتند از:
\begin{itemize}
    \item پیکربندی ثبّات‌های نگاشت‌شده در حافظه برای واحدهای \lr{ETMv4}، \lr{CSTF} و \lr{TPIU} با استفاده از اسکریپت‌های کنترل دیباگر \lr{GDB} و \lr{Firmware} سیستم.
    \item استخراج فیزیکی سیگنال‌های ۵ سیمه \lr{Trace} بر روی بستر سخت‌افزاری دست‌ساز با حفظ سلامت سیگنال.
    \item توسعه مبدل نرم‌افزاری پیشرفته جهت تبدیل داده‌های ناهمگام و خام لاجیک‌آنالایزر به فریم‌های ساختاریافته \lr{OpenCSD}.
    \item بازسازی پیوسته بیش از دویست هزار دستور اجرایی با نرخ خطای صفر و استخراج تفکیک‌شده مدت زمان توابع و وقفه‌ها.
    \item اثبات عملیاتی بودن یک روش کاملاً اقتصادی و مبتنی بر سخت‌افزارهای تجاری در دسترس به عنوان جایگزین پروب‌های فوق‌گران‌قیمت صنعتی.
\end{itemize}

\section{نوآوری‌ها و دستاوردهای شاخص}
دستاوردهای کلیدی حاصل از این پژوهش به شرح زیر خلاصه می‌شود:
\begin{enumerate}
    \item \textbf{حذف کامل اثر پروب و دستیابی به جیتر صفر (\lr{Zero Jitter}):} اثبات عملیاتی این حقیقت که ردگیری سخت‌افزاری کدهای برنامه، حتی در تکرار ۴۹۲ بار اجرای روتین وقفه \lr{SysTick}، همواره طول اجرای دقیقاً یکسان (۸ دستورالعمل) را ثبت می‌کند که نشانه انطباق مطلق با الزامات سامانه‌های بی‌درنگ سخت است.
    \item \textbf{حل چالش بحرانی بسته‌های همگام‌سازی (\lr{A-Sync Issue}):} ابداع رویکرد سوئیچینگ کنترل‌شده بیت فعال‌سازی \lr{ETM} (\lr{TRCPRGCTLR.EN}) در مبادی پنجره‌های \lr{Trace} کوچک، جهت تولید اجباری بسته‌های همگام‌سازی اولیه و نجات جریان داده از حالت غیرقابل‌رمزگشایی (\lr{NO\_SYNC}).
    \item \textbf{بومی‌سازی و ساخت سخت‌افزار سفارشی در شرایط کمبود قطعات:} طراحی و مونتاژ موفق هدربرد اختصاصی \lr{STM32H750} در بازه زمانی فشرده یک هفته‌ای.
    \item \textbf{طراحی خط لوله نرم‌افزاری خودکار و اتصال به \lr{OpenCSD}:} ایجاد پیوند میان خروجی‌های خام لاجیک‌آنالایزرهای اقتصادی و کتابخانه مرجع صنعتی رمزگشایی \lr{Trace}.
\end{enumerate}

\section{ساختار فصول پایان‌نامه}
این رساله در هشت فصل مدون شده است:
\begin{itemize}\setlength{\itemsep}{0.4em}
    \item \textbf{فصل دوم (مفاهیم پایه و پیشینه پژوهش):} کالبدشکافی تئوری‌های \lr{WCET}، ساختار زیرسیستم عیب‌یابی و ردگیری \lr{ARM CoreSight} و مقایسه تطبیقی رویکردهای صنعتی.
    \item \textbf{فصل سوم (معماری کلان سامانه):} تشریح مدل مفهومی سامانه، لایه‌های شش‌گانه جریان داده و نحوه تعامل با سامانه آزمون خودکار.
    \item \textbf{فصل چهارم (طراحی و پیاده‌سازی سخت‌افزار):} مشخصات برد سفارشی، اتصالات و سیم‌کشی پایه‌ها، ملاحظات فرکانس بالا و راهبرد نمونه‌برداری.
    \item \textbf{فصل پنجم (پیکربندی \lr{ETMv4} و \lr{Firmware}):} پروتکل فشرده‌سازی \lr{ETMv4}، راه‌اندازی ثبّاتی از طریق اسکریپت‌های \lr{GDB} و حل چالش \lr{A-Sync}.
    \item \textbf{فصل ششم (پردازش و رمزگشایی \lr{Trace}):} معماری اسکریپت \lr{TraceStreamProcessor}، ساختاربندی اسنپ‌شات‌ها و رمزگشایی دستور به دستور با \lr{OpenCSD}.
    \item \textbf{فصل هفتم (ارزیابی تجربی و نتایج):} تحلیل عمیق نتایج آزمون‌های \lr{TimerTest}، \lr{Routine\_Irq} و \lr{FreqTest} و ارزیابی اقتصادی و مقایسه‌ای بسترها.
    \item \textbf{فصل هشتم (نتیجه‌گیری و کارهای آینده):} مرور دستاوردها، درس‌آموخته‌های مهندسی، جمع‌بندی اقتصادی پلتفرم و چشم‌انداز پیاده‌سازی \lr{FPGA}.
\end{itemize}
